% GENERATED FILE. Edit canonical lessons, then run npm run generate:books.
% canonical-source-hash: fnv1a64:a066c03e07d40048
% canonical-lessons: TE-C502-nitturcu, TE-C502-nirasapadu, TE-C502-nirutsaham, TE-C502-tahatahaladu, TE-C502-krsincu

\chapter{Sighs and Letdowns}
\label{ch:te-sighs-and-letdowns}
\hlchaptermodality{502}

\begin{chapteropening}
\textbf{By the end of this chapter:} \emph{I can say that I sighed, was disappointed or longed for something.}

Complete the last lesson of chapter 502: I can say that I sighed, was disappointed or longed for something.
\end{chapteropening}

\section[\texorpdfstring{\te{నిట్టూర్చు} (niṭṭūrcu)}{నిట్టూర్చు (niṭṭūrcu)}]{\te{నిట్టూర్చు} (niṭṭūrcu) — to sigh}

\label{lesson:TE-C502-nitturcu}



\begin{quote}
Before the new one: say the Telugu for to impress, then the Telugu for to cheer on.
\end{quote}

\subsection*{You'll want to know: \te{నిట్టూర్చు}}
\textbf{\te{నిట్టూర్చు}} — \emph{niṭṭūrcu} — \textquotedblleft{}to sigh\textquotedblright{}.

\begin{quote}
\te{ఫలితం} \te{చూసి} \te{అతను} \te{నిట్టూర్చాడు}.
\end{quote}

\begin{quote}
\emph{phalitaṁ cūsi atanu niṭṭūrcāḍu.}
\end{quote}

\begin{quote}
Seeing the result, he sighed.
\end{quote}

Past (I …): \te{నిట్టూర్చాను}, \emph{niṭṭūrcānu}.

The noun is \te{నిట్టూర్పు}, a sigh.

\begin{cousinweb}
\textbf{\te{నిట్టూర్చు}} (\emph{niṭṭūrcu}) is native Dravidian: DEDR 3738 lists Telugu \emph{niṭṭ-ūrucu}, \textquotedblleft{}to sigh\textquotedblright{}, and \emph{niṭṭ-ūrupu}, \textquotedblleft{}a sigh\textquotedblright{}, in the article of Tamil \emph{neṭu}, \textquotedblleft{}long\textquotedblright{}: a long breath.
\end{cousinweb}

\begin{grammarlens}[title={What you've built}]
One more word for describing what you lived and felt.
\end{grammarlens}

\subsection*{Your turn}
\begin{itemize}
\raggedright
  \item \emph{Say it:} \emph{niṭṭūrcu}
  \item \emph{Say it:} \emph{niṭṭūrcu}, once more
  \item \emph{Say it:} \emph{phalitaṁ cūsi atanu niṭṭūrcāḍu.} — Seeing the result, he sighed.
  \item \emph{Recall:} say the Telugu for to impress, then the Telugu for to cheer on, then say \emph{niṭṭūrcu} again
\end{itemize}

\begin{tcolorbox}[breakable,colback=teal!4,colframe=teal!35!black,title={Before you move on}]
What does \te{నిట్టూర్చు} mean? (\textquotedblleft{}To sigh\textquotedblright{}.) Say it once more.
\end{tcolorbox}
\section[\texorpdfstring{\te{నిరాశపడు} (nirāśapaḍu)}{నిరాశపడు (nirāśapaḍu)}]{\te{నిరాశపడు} (nirāśapaḍu) — to be disappointed}

\label{lesson:TE-C502-nirasapadu}



\begin{quote}
Before the new one: say the Telugu for to cheer on, then the Telugu for to sigh.
\end{quote}

\subsection*{You'll want to know: \te{నిరాశపడు}}
\textbf{\te{నిరాశపడు}} — \emph{nirāśapaḍu} — \textquotedblleft{}to be disappointed\textquotedblright{}.

\begin{quote}
\te{వర్షం} \te{వల్ల} \te{ఆట} \te{రద్దు} \te{అయింది}, \te{పిల్లలు} \te{నిరాశపడ్డారు}.
\end{quote}

\begin{quote}
\emph{varṣaṁ valla āṭa raddu ayindi, pillalu nirāśapaḍḍāru.}
\end{quote}

\begin{quote}
The game was called off because of rain; the children were disappointed.
\end{quote}

Past (I …): \te{నిరాశపడ్డాను}, \emph{nirāśapaḍḍānu}.

\te{నిరాశ}, disappointment, + \te{పడు}. \te{నిరాశపరచు}, to disappoint someone.

\begin{cousinweb}
\textbf{\te{నిరాశపడు}} (\emph{nirāśapaḍu}) is built from \te{నిరాశ}, from Sanskrit \emph{nirāśa}, \textquotedblleft{}without hope, despondent\textquotedblright{} (Monier-Williams) and the native verb \emph{paḍu}, \textquotedblleft{}to feel, to suffer\textquotedblright{} (DEDR 3853).
\end{cousinweb}

\begin{grammarlens}[title={What you've built}]
One more word for describing what you lived and felt.
\end{grammarlens}

\subsection*{Your turn}
\begin{itemize}
\raggedright
  \item \emph{Say it:} \emph{nirāśapaḍu}
  \item \emph{Say it:} \emph{nirāśapaḍu}, once more
  \item \emph{Say it:} \emph{varṣaṁ valla āṭa raddu ayindi, pillalu nirāśapaḍḍāru.} — The game was called off because of rain; the children were disappointed.
  \item \emph{Recall:} say the Telugu for to cheer on, then the Telugu for to sigh, then say \emph{nirāśapaḍu} again
\end{itemize}

\begin{tcolorbox}[breakable,colback=teal!4,colframe=teal!35!black,title={Before you move on}]
What does \te{నిరాశపడు} mean? (\textquotedblleft{}To be disappointed\textquotedblright{}.) Say it once more.
\end{tcolorbox}
\section[\texorpdfstring{\te{నిరుత్సాహం} (nirutsāhaṁ)}{నిరుత్సాహం (nirutsāhaṁ)}]{\te{నిరుత్సాహం} (nirutsāhaṁ) — discouragement}

\label{lesson:TE-C502-nirutsaham}



\begin{quote}
Before the new one: say the Telugu for to sigh, then the Telugu for to be disappointed.
\end{quote}

\subsection*{You'll want to know: \te{నిరుత్సాహం}}
\textbf{\te{నిరుత్సాహం}} — \emph{nirutsāhaṁ} — \textquotedblleft{}discouragement\textquotedblright{}.

\begin{quote}
\te{మొదటిసారి} \te{గెలవకపోయినా} \te{ఆమెకు} \te{నిరుత్సాహం} \te{లేదు}.
\end{quote}

\begin{quote}
\emph{modaṭisāri gelavakapōyinā āmeku nirutsāhaṁ lēdu.}
\end{quote}

\begin{quote}
Though she did not win the first time, she was not discouraged.
\end{quote}

The opposite of \te{ఉత్సాహం}, enthusiasm: \te{నిర్}-, without.

\begin{cousinweb}
\textbf{\te{నిరుత్సాహం}} (\emph{nirutsāhaṁ}) is borrowed from Sanskrit \emph{nirutsāha}, \textquotedblleft{}lacking energy or heart\textquotedblright{} (Monier-Williams): \emph{nir-}, \textquotedblleft{}without\textquotedblright{}, and \emph{utsāha}, \textquotedblleft{}power, firmness\textquotedblright{} (CDIAL 1882).
\end{cousinweb}

\begin{grammarlens}[title={What you've built}]
One more word for describing what you lived and felt.
\end{grammarlens}

\subsection*{Your turn}
\begin{itemize}
\raggedright
  \item \emph{Say it:} \emph{nirutsāhaṁ}
  \item \emph{Say it:} \emph{nirutsāhaṁ}, once more
  \item \emph{Say it:} \emph{modaṭisāri gelavakapōyinā āmeku nirutsāhaṁ lēdu.} — Though she did not win the first time, she was not discouraged.
  \item \emph{Recall:} say the Telugu for to sigh, then the Telugu for to be disappointed, then say \emph{nirutsāhaṁ} again
\end{itemize}

\begin{tcolorbox}[breakable,colback=teal!4,colframe=teal!35!black,title={Before you move on}]
What does \te{నిరుత్సాహం} mean? (\textquotedblleft{}Discouragement\textquotedblright{}.) Say it once more.
\end{tcolorbox}
\section[\texorpdfstring{\te{తహతహలాడు} (tahatahalāḍu)}{తహతహలాడు (tahatahalāḍu)}]{\te{తహతహలాడు} (tahatahalāḍu) — to yearn; to be restless for}

\label{lesson:TE-C502-tahatahaladu}



\begin{quote}
Before the new one: say the Telugu for to be disappointed, then the Telugu for discouragement.
\end{quote}

\subsection*{You'll want to know: \te{తహతహలాడు}}
\textbf{\te{తహతహలాడు}} — \emph{tahatahalāḍu} — \textquotedblleft{}to yearn; to be restless for\textquotedblright{}.

\begin{quote}
\te{అమ్మను} \te{చూడాలని} \te{పాప} \te{తహతహలాడింది}.
\end{quote}

\begin{quote}
\emph{ammanu cūḍālani pāpa tahatahalāḍindi.}
\end{quote}

\begin{quote}
The little girl longed to see her mum.
\end{quote}

Past (I …): \te{తహతహలాడాను}, \emph{tahatahalāḍānu}.

An echo word, \te{తహతహ}, + \te{ఆడు}. -\te{ఆలని}, wanting to.

\begin{cousinweb}
\textbf{\te{తహతహలాడు}} (\emph{tahatahalāḍu}) is built from the doubled echo word \te{తహతహ}, \textquotedblleft{}restless longing\textquotedblright{}, and the verb \emph{āḍu}, \textquotedblleft{}to move\textquotedblright{} (DEDR 347). The origin of \te{తహతహ} is not settled.
\end{cousinweb}

\begin{grammarlens}[title={What you've built}]
One more word for describing what you lived and felt.
\end{grammarlens}

\subsection*{Your turn}
\begin{itemize}
\raggedright
  \item \emph{Say it:} \emph{tahatahalāḍu}
  \item \emph{Say it:} \emph{tahatahalāḍu}, once more
  \item \emph{Say it:} \emph{ammanu cūḍālani pāpa tahatahalāḍindi.} — The little girl longed to see her mum.
  \item \emph{Recall:} say the Telugu for to be disappointed, then the Telugu for discouragement, then say \emph{tahatahalāḍu} again
\end{itemize}

\begin{tcolorbox}[breakable,colback=teal!4,colframe=teal!35!black,title={Before you move on}]
What does \te{తహతహలాడు} mean? (\textquotedblleft{}To yearn; to be restless for\textquotedblright{}.) Say it once more.
\end{tcolorbox}
\section[\texorpdfstring{\te{కృశించు} (kṛśiñcu)}{కృశించు (kṛśiñcu)}]{\te{కృశించు} (kṛśiñcu) — to waste away; to pine}

\label{lesson:TE-C502-krsincu}



\begin{quote}
Before the new one: say the Telugu for discouragement, then the Telugu for to yearn.
\end{quote}

\subsection*{You'll want to know: \te{కృశించు}}
\textbf{\te{కృశించు}} — \emph{kṛśiñcu} — \textquotedblleft{}to waste away; to pine\textquotedblright{}.

\begin{quote}
\te{దిగులుతో} \te{ఆమె} \te{రోజురోజుకీ} \te{కృశించింది}.
\end{quote}

\begin{quote}
\emph{digulutō āme rōjurōjukī kṛśiñcindi.}
\end{quote}

\begin{quote}
From grief she wasted away day by day.
\end{quote}

Past (I …): \te{కృశించాను}, \emph{kṛśiñcānu}.

A formal word for growing thin from illness or sorrow.

\begin{cousinweb}
\textbf{\te{కృశించు}} (\emph{kṛśiñcu}) is built on Sanskrit \emph{kṛśa}, \textquotedblleft{}lean, emaciated, weak\textquotedblright{} (Monier-Williams), with the Telugu verb ending -\te{ఇంచు}.
\end{cousinweb}

\begin{grammarlens}[title={What you've built}]
One more word for describing what you lived and felt.
\end{grammarlens}

\subsection*{Your turn}
\begin{itemize}
\raggedright
  \item \emph{Say it:} \emph{kṛśiñcu}
  \item \emph{Say it:} \emph{kṛśiñcu}, once more
  \item \emph{Say it:} \emph{digulutō āme rōjurōjukī kṛśiñcindi.} — From grief she wasted away day by day.
  \item \emph{Recall:} say the Telugu for discouragement, then the Telugu for to yearn, then say \emph{kṛśiñcu} again
\end{itemize}

\begin{tcolorbox}[breakable,colback=teal!4,colframe=teal!35!black,title={Before you move on}]
What does \te{కృశించు} mean? (\textquotedblleft{}To waste away; to pine\textquotedblright{}.) Say it once more.
\end{tcolorbox}
